\section{Filosofía de la ingeniería y metodología}

\subsection{Enseñarle Engineering a un Researcher es más difícil que al revés}

Weng Jiayi (翁家翌) cita a un colega (un PhD que trabajó en un RL framework reconocido):

\begin{importantbox}{La asimetría entre Engineering y Research}
«Enseñarle a un researcher a hacer bien engineering es mucho más difícil que enseñarle a un engineer a hacer bien research.» En la era actual de la IA, el valor de la capacidad de ingeniería supera al de la capacidad académica, porque la intuición de research puede acumularse trabajando mucho tiempo en la industria, mientras que una capacidad de ingeniería sólida requiere un entrenamiento deliberado y prolongado.
\end{importantbox}

\subsection{La velocidad de iteración lo determina todo}

La idea central que Weng Jiayi (翁家翌) enfatiza una y otra vez es que \textbf{el número de iteraciones correctas por unidad de tiempo} es el factor determinante del éxito en la investigación de IA. En concreto:

\begin{enumerate}
    \item Corrección de la infra: una infra sin bugs significa que el resultado de cada experimento es confiable
    \item Velocidad de iteración: pasar de 30 iteraciones por semana a 300 iteraciones por semana hace que la tasa de éxito crezca de forma lineal
    \item Rendimiento en la corrección de bugs: cuántos bugs se pueden corregir por unidad de tiempo, cuántas iteraciones correctas se pueden lograr
\end{enumerate}

\begin{warningbox}{La innovación algorítmica no es el cuello de botella}
Weng Jiayi (翁家翌) considera que el cuello de botella de la frontera actual de la IA \textbf{no está en la innovación algorítmica}. «Si corriges todos los bugs, es posible que el algoritmo funcione bien sin necesidad de cambiar nada.» Muchos problemas que parecen requerir un nuevo algoritmo para resolverse en realidad pueden deberse solo a bugs de la infra. Esto contrasta claramente con el paradigma académico, centrado en proponer nuevos algoritmos.
\end{warningbox}

\subsection{Consistencia del código y deterioro de los proyectos}

Weng Jiayi (翁家翌) considera la \textbf{consistencia (consistency)} como el primer principio tanto de los proyectos de software como de la gestión de organizaciones. La raíz del deterioro de los proyectos es la falta de consistencia: cuando varias personas colaboran, el context de cada una es distinto, los supuestos no pueden transmitirse, y esto provoca que el código se infle y su calidad se deteriore.

Él considera que dirigir una empresa y mantener un código base son esencialmente lo mismo: ambos requieren mantener la consistency. Si no hay consistencia, el código base se vuelve inflado y la estructura organizacional también se vuelve inflada. La solución es que \textbf{la información fluya sin obstáculos}: que las decisiones de arriba lleguen abajo sin pérdida, y que el progreso de abajo llegue arriba sin pérdida.

Weng Jiayi (翁家翌) lleva este planteamiento hasta el extremo: en teoría, el \textbf{context sharing debería ser reemplazado por un agent con un context de longitud infinita}. Esto porque el context del cerebro humano es limitado y no puede almacenar al mismo tiempo toda la información de una organización entera. Pero la IA sí puede. En el futuro, tal vez cada empresa tenga un agent de context infinito de este tipo que funja como CEO, encargado de todo el sharing y las decision: «Probablemente no haya un decision maker más adecuado que un agent así.» Esta perspectiva unifica su filosofía de la ingeniería, sus convicciones sobre la IA y sus reflexiones sobre la gestión organizacional.

\subsection{Resumen de la sección}

La filosofía de la ingeniería de Weng Jiayi (翁家翌) puede resumirse en tres principios centrales: (1) \textbf{prioridad a la consistencia}: que una sola persona se encargue de todo, lo cual, aunque no favorece la escalabilidad a largo plazo, garantiza la consistencia; (2) \textbf{prioridad a la velocidad de iteración}: no se busca un algoritmo genial de una sola vez, sino aproximarse a la solución óptima mediante prueba y error rápidos; (3) \textbf{prioridad a la corrección}: corregir bugs es más importante que escribir un nuevo feature.


